% !TEX root = ../main.tex
\corrigesChapitre{6}{Schémas de Lewis, géométrie et polarité}
\corrigesSouples

\begin{blocLivre}{cor-6-1}
\begin{corrige}{1}{Électrons de valence et liaisons}
\demarche{Les électrons de valence sont ceux de la dernière couche occupée ; il en manque autant que de liaisons à former.}
\begin{questions}
\item H : $1\mathrm{s}^1$ ; C : $1\mathrm{s}^2\,2\mathrm{s}^2\,2\mathrm{p}^2$ ; N : $1\mathrm{s}^2\,2\mathrm{s}^2\,2\mathrm{p}^3$ ; O : $1\mathrm{s}^2\,2\mathrm{s}^2\,2\mathrm{p}^4$ ;
Cl : $1\mathrm{s}^2\,2\mathrm{s}^2\,2\mathrm{p}^6\,3\mathrm{s}^2\,3\mathrm{p}^5$. Électrons de valence : \res{\text{1, 4, 5, 6 et 7}}.
\item Il en manque respectivement \res{\text{1, 4, 3, 2 et 1}}.
\item H forme \res{1} liaison, C \res{4}, N \res{3}, O \res{2} et Cl \res{1}.
\item Doublets non liants : H \res{0}, C \res{0}, N \res{1}, O \res{2}, Cl \res{3}.
\end{questions}
\retiens{Le nombre de liaisons qu'un atome forme est le nombre d'électrons qui lui manquent pour atteindre le duet ou
l'octet ; ses autres électrons de valence forment ses doublets non liants.}
\end{corrige}
\end{blocLivre}
\begin{blocLivre}{cor-6-2}
\begin{corrige}{2}{Vérifier des schémas de Lewis}
\demarche{Une liaison double compte deux doublets liants, une liaison triple trois ; le total des doublets vaut la moitié des
électrons de valence.}
\begin{questions}
\item Méthanol : \res{5} liants et \res{2} non liants ; éthène : \res{6} liants (dont deux pour la liaison double), \res{0} non liant ;
cyanure d'hydrogène : \res{4} liants et \res{1} non liant ; peroxyde d'hydrogène : \res{3} liants et \res{4} non liants.
\item Méthanol : $4+6+4\times1=14$ électrons, 7 doublets ; éthène : $2\times4+4\times1=12$, 6 doublets ; $\ce{HCN}$ : $1+4+5=10$, 5 doublets ;
$\ce{H2O2}$ : $2\times1+2\times6=14$, 7 doublets. Les schémas sont \res{\text{cohérents}}.
\item Chaque atome d'hydrogène est entouré de 2 électrons, chaque atome C, N ou O de 8 : les règles sont
\res{\text{respectées}}.
\item Liaison double : l'\res{\text{éthène}} ; liaison triple : le \res{\text{cyanure d'hydrogène}}.
\end{questions}
\end{corrige}
\end{blocLivre}
\begin{blocLivre}{cor-6-3}
\begin{corrige}{3}{Des ions}
\demarche{Pour un ion positif, on retranche des électrons ; pour un ion négatif, on en ajoute.}
\begin{questions}
\item \chemfig{H-\charge{90=\|,45:3pt=\cplus}{O}(-[6]H)-H} \quad \chemfig{\charge{90=\|,180=\|,270=\|,135:3pt=\cmoins}{O}-H} \quad
\chemfig{H-\charge{45:3pt=\cplus}{N}(-[2]H)(-[6]H)-H} \quad \chemfig{\charge{0=\|,90=\|,180=\|,270=\|,45:3pt=\cmoins}{Cl}}
\item L'ion $\ce{H+}$ ne possède aucun électron : il présente une \res{\text{lacune électronique}}.
\item Un \res{\text{doublet non liant de l'oxygène}} de la molécule d'eau vient combler la lacune de l'ion $\ce{H+}$ : il devient un
doublet liant.
\item L'oxygène partage désormais ce doublet qui lui appartenait en propre : il « perd » l'équivalent d'un électron, d'où la
charge \res{\oplus}.
\end{questions}
\end{corrige}
\end{blocLivre}
\begin{blocLivre}{cor-6-4}
\begin{corrige}{4}{Prévoir une géométrie}
\demarche{On compte les directions de doublets autour de l'atome central ; une liaison multiple compte pour une direction.}
\begin{questions}
\item $\ce{CH4}$ : 4 liants, 0 non liant ; $\ce{NH3}$ : 3 liants, 1 non liant ; $\ce{H2O}$ : 2 liants, 2 non liants.
\item \res{\text{Tétraédrique}}, \res{\text{pyramidale}} et \res{\text{coudée}}.
\item Deux directions dans les deux cas : géométrie \res{\text{linéaire}}.
\item Trois directions autour du carbone : molécule \res{\text{plane}}, d'angles voisins de \qty{120}{\degree}.
\end{questions}
\retiens{On compte les directions de doublets autour de l'atome central, doublets non liants compris ; mais on nomme la
géométrie d'après la seule position des atomes.}
\end{corrige}
\end{blocLivre}
\begin{blocLivre}{cor-6-5}
\begin{corrige}{5}{Liaisons polarisées}
\demarche{Une liaison est polarisée si l'écart d'électronégativité atteint 0,4 ; $\delta^{-}$ va sur l'atome le plus électronégatif.}
\begin{questions}
\item H--Cl : 0,96 ; C--H : 0,35 ; O--H : 1,24 ; C--O : 0,89 ; C--Cl : 0,61 ; N--H : 0,84 ; H--H et C--C : 0.
\item Liaisons polarisées : \res{\text{H--Cl, O--H, C--O, C--Cl et N--H}}.
\item $\mathrm{H}^{\delta+}$--$\mathrm{Cl}^{\delta-}$ ; $\mathrm{O}^{\delta-}$--$\mathrm{H}^{\delta+}$ ; $\mathrm{C}^{\delta+}$--$\mathrm{O}^{\delta-}$ ; $\mathrm{C}^{\delta+}$--$\mathrm{Cl}^{\delta-}$ ; $\mathrm{N}^{\delta-}$--$\mathrm{H}^{\delta+}$.
\item La liaison \res{\text{O--H}} est la plus polarisée. Entre deux atomes identiques, l'écart est nul : la liaison n'est
\res{\text{pas polarisée}}.
\end{questions}
\end{corrige}
\end{blocLivre}
\begin{blocLivre}{cor-6-6}
\begin{corrige}{6}{Compter les doublets en Python}
\demarche{Le programme compte les électrons de valence, retranche la charge, puis prend le quotient entier par 2.}
\begin{questions}
\item Le programme affiche \res{4}, \res{4}, \res{4} et \res{4} : ces quatre entités comptent chacune quatre doublets.
\item Un ion positif a perdu des électrons, un ion négatif en a gagné : retrancher la charge corrige le compte.
\code{//} donne le \res{\text{quotient entier}} de la division.
\item $\ce{HCN}$ : $\dfrac{1+4+5}{2}=\res{5}$ ; $\ce{C2H6O}$ : $\dfrac{2\times4+6\times1+6}{2}=\res{10}$.
\item 17 est impair : $\dfrac{17}{2}$ donne 8 doublets et il reste un \res{\text{électron célibataire}}. La division entière le fait
disparaître ; une telle molécule, appelée radical, ne peut pas respecter l'octet pour tous ses atomes.
\end{questions}
\end{corrige}
\end{blocLivre}
\begin{blocLivre}{cor-6-7}
\begin{corrige}{7}{Le méthanal}
\demarche{$4+2\times1+6=12$ électrons de valence, soit 6 doublets : 4 liants, dont une liaison double, et 2 non liants sur l'oxygène.}
\begin{questions}
\item \res{6 \text{ doublets}} : \chemfig{\charge{135=\|,225=\|}{O}=C(-[1]H)-[7]H}
\item \res{\text{Trois directions}} autour du carbone : la molécule est \res{\text{plane}}, avec des angles voisins de \qty{120}{\degree}.
\item C=O est polarisée ($0{,}89$) : $\delta^{+}$ sur C, $\delta^{-}$ sur O ; les liaisons C--H ne le sont pas.
\item G$^{+}$ est sur le carbone, G$^{-}$ sur l'oxygène : la molécule est \res{\text{polaire}}.
\end{questions}
\end{corrige}
\end{blocLivre}
\begin{blocLivre}{cor-6-8}
\begin{corrige}{8}{Un filet de liquide dévié}
\demarche{Seules les molécules polaires sont attirées par un objet chargé : elles s'orientent, pôle de signe opposé vers la règle.}
\begin{questions}
\item Le filet d'\res{\text{eau}} est dévié vers la règle ; le filet de cyclohexane ne l'est pas.
\item Liaisons O--H polarisées ($1{,}24$) et géométrie coudée : G$^{+}$, entre les deux hydrogènes, est distinct de G$^{-}$, sur
l'oxygène. La molécule est \res{\text{polaire}}.
\item Les liaisons C--C (écart nul) et C--H ($0{,}35$) ne sont \res{\text{pas polarisées}} : la molécule ne porte pas de charges
partielles.
\item Les molécules d'eau tournent leur pôle $\delta^{-}$ vers la règle positive et sont \res{\text{attirées}}. Avec une règle négative,
elles tournent leur pôle $\delta^{+}$ vers elle : le filet est encore \res{\text{attiré}}.
\end{questions}
\schema[corrige]{06-ex08-cor}
\end{corrige}
\end{blocLivre}
\begin{blocLivre}{cor-6-9}
\begin{corrige}{9}{L'ammoniac et l'ion ammonium}
\demarche{L'ion ammonium compte $\dfrac{5+4\times1-1}{2}=4$ doublets, tous liants ; l'ammoniac a 3 doublets liants et 1 non liant.}
\begin{questions}
\item \chemfig{H-\charge{90=\|}{N}(-[6]H)-H} \qquad \chemfig{H-\charge{45:3pt=\cplus}{N}(-[2]H)(-[6]H)-H}
\item Ammoniac : \res{\text{pyramidale}} ; ion ammonium : \res{\text{tétraédrique}}.
\item Le doublet non liant de l'ammoniac, plus proche du noyau d'azote, \res{\text{repousse davantage}} les doublets liants : l'angle
est plus petit que \qty{109,5}{\degree}. Dans l'ion ammonium, les quatre doublets sont identiques.
\item Le \res{\text{doublet non liant}} de l'azote comble la \res{\text{lacune}} de l'ion $\ce{H+}$ et forme une quatrième liaison N--H ; l'azote
porte alors la charge $\oplus$.
\end{questions}
\end{corrige}
\end{blocLivre}
\begin{blocLivre}{cor-6-10}
\begin{corrige}{10}{Deux solvants chlorés}
\demarche{$\ce{CHCl3}$ : $\dfrac{4+1+3\times7}{2}=13$ doublets ; $\ce{CCl4}$ : $\dfrac{4+4\times7}{2}=16$ doublets ; chaque chlore porte 3 doublets non liants.}
\begin{questions}
\item \chemfig{H-C(-[2]\charge{0=\|,90=\|,180=\|}{Cl})(-[6]\charge{0=\|,180=\|,270=\|}{Cl})-\charge{0=\|,90=\|,270=\|}{Cl}} \qquad
\chemfig{\charge{90=\|,180=\|,270=\|}{Cl}-C(-[2]\charge{0=\|,90=\|,180=\|}{Cl})(-[6]\charge{0=\|,180=\|,270=\|}{Cl})-\charge{0=\|,90=\|,270=\|}{Cl}}
\item Quatre doublets liants autour du carbone : géométrie \res{\text{tétraédrique}} pour les deux.
\item Les liaisons \res{\text{C--Cl}} ($0{,}61$) sont polarisées, pas la liaison C--H.
\item Dans $\ce{CCl4}$, les quatre liaisons identiques sont disposées symétriquement : G$^{-}$ est au centre du tétraèdre, sur le
carbone, comme G$^{+}$ : molécule \res{\text{apolaire}}. Dans $\ce{CHCl3}$, G$^{-}$, au centre des trois chlores, est distinct de G$^{+}$ :
molécule \res{\text{polaire}}.
\end{questions}
\schema[corrige]{06-ex10-cor}
\piege{Conclure qu'une molécule est polaire dès qu'elle contient des liaisons polarisées : si ces liaisons sont disposées
symétriquement, leurs effets se compensent et la molécule est apolaire.}
\end{corrige}
\end{blocLivre}
\begin{blocLivre}{cor-6-11}
\begin{corrige}{11}{Deux molécules, une même formule}
\demarche{$2\times4+6\times1+6=20$ électrons de valence : 10 doublets, dont 8 liants et 2 non liants sur l'oxygène.}
\begin{questions}
\item \res{10 \text{ doublets}} pour chacune.
\item \chemfig{H-C(-[2]H)(-[6]H)-C(-[2]H)(-[6]H)-\charge{90=\|,270=\|}{O}-H} \qquad \chemfig{H-C(-[2]H)(-[6]H)-\charge{90=\|,270=\|}{O}-C(-[2]H)(-[6]H)-H}
\item Autour de l'oxygène : deux liants et deux non liants, géométrie \res{\text{coudée}} ; autour de chaque carbone,
\res{\text{tétraédrique}}.
\item Les liaisons C--O et O--H sont polarisées et la géométrie est coudée autour de l'oxygène : les deux molécules sont
\res{\text{polaires}}.
\end{questions}
\end{corrige}
\end{blocLivre}
\begin{blocLivre}{cor-6-12}
\begin{corrige}{12}{L'eau dans le four à micro-ondes}
\demarche{Une molécule polaire s'oriente dans un champ électrique ; s'il change de sens sans cesse, elle tourne et s'agite.}
\begin{questions}
\item \chemfig{H-\charge{90=\|,270=\|}{O}-H} : 2 doublets liants et 2 non liants.
\item Quatre directions de doublets autour de l'oxygène, dont deux non liantes : la molécule est \res{\text{coudée}}. Les doublets non
liants repoussent plus fortement les doublets liants : l'angle est \res{\text{plus petit}} que \qty{109,5}{\degree}.
\item $3{,}44-2{,}20=1{,}24$ : liaisons polarisées ; G$^{+}$, au milieu des deux hydrogènes, est distinct de G$^{-}$, sur l'oxygène :
la molécule est \res{\text{polaire}}.
\item L'eau, polaire, tend à s'orienter dans le champ électrique ; les molécules apolaires du liquide modèle ne possèdent pas ce
dipôle permanent. Cela ne suffit pas à prévoir une vitesse de chauffage : l'absorption, la masse et la capacité thermique du
liquide interviennent aussi. Une huile alimentaire réelle ne se réduit pas à ce liquide modèle.
\end{questions}
\retiens{Pour conclure sur la polarité, il faut à la fois des liaisons polarisées et une géométrie non symétrique.}
\end{corrige}
\end{blocLivre}
\begin{blocLivre}{cor-6-13}
\begin{corrige}{13}{Les fluides frigorigènes}
\demarche{$\ce{CCl2F2}$ : $\dfrac{4+2\times7+2\times7}{2}=16$ doublets ; $\ce{CH2F2}$ : $\dfrac{4+2+2\times7}{2}=10$ doublets ; les halogènes portent 3 doublets non liants.}
\begin{questions}
\item \res{16} et \res{10} doublets : 4 liants autour du carbone et 3 non liants sur chaque atome de chlore ou de fluor.
\par\smallskip\chemfig{\charge{90=\|,180=\|,270=\|}{F}-C(-[2]\charge{0=\|,90=\|,180=\|}{Cl})(-[6]\charge{0=\|,180=\|,270=\|}{Cl})-\charge{0=\|,90=\|,270=\|}{F}}
\qquad\chemfig{\charge{90=\|,180=\|,270=\|}{F}-C(-[2]H)(-[6]H)-\charge{0=\|,90=\|,270=\|}{F}}
\item Quatre directions autour du carbone : \res{\text{tétraédrique}} pour les deux.
\item C--Cl ($0{,}61$) et C--F ($1{,}43$) sont polarisées ; C--H ($0{,}35$) ne l'est pas.
\item Dans $\ce{CCl2F2}$ et $\ce{CH2F2}$, les liaisons polarisées ne sont pas toutes identiques : G$^{-}$ n'est pas sur le carbone, ces
molécules sont polaires. Le \res{\text{dioxyde de carbone}}, linéaire et symétrique, est le seul fluide apolaire.
\end{questions}
\end{corrige}
\end{blocLivre}
\begin{blocLivre}{cor-6-14}
\begin{corrige}{14}{Un gaz incolore et mortel}
\demarche{Deux entités qui ont le même nombre d'électrons de valence peuvent avoir des schémas de Lewis semblables.}
\begin{questions}
\item \chemfig{\charge{180=\|}{N}~\charge{0=\|}{N}} : chaque azote est entouré de $3\times2+2=8$ électrons, l'octet est \res{\text{respecté}}.
\item $\ce{N2}$ : $2\times5=10$ ; $\ce{CO}$ : $4+6=10$ : \res{\text{même nombre}}, soit 5 doublets.
\item \chemfig{\charge{180=\|,135:4pt=\cmoins}{C}~\charge{0=\|,45:4pt=\cplus}{O}} : une liaison triple et un doublet non liant sur chaque atome. Le carbone ne
forme ici que \res{3 \text{ liaisons}}, au lieu de 4 habituellement. Les charges formelles sont $-1$ sur C et $+1$ sur O ; leur somme est nulle.
\item Une molécule diatomique est \res{\text{linéaire}}. D'après la mesure fournie, les centres des charges sont distincts : CO est
\res{\text{polaire}}. Son faible dipôle expérimental ne se déduit pas de la seule différence d'électronégativité.
\end{questions}
\piege{Les charges formelles du schéma de Lewis servent au décompte des électrons ; elles ne sont pas des charges partielles mesurées. La molécule CO est globalement neutre.}
\end{corrige}
\end{blocLivre}
\begin{blocLivre}{cor-6-15}
\begin{corrige}{15}{Identifier des modèles moléculaires}
\demarche{Les tiges donnent les liaisons et leur multiplicité ; on vérifie ensuite que chaque atome respecte le duet ou l'octet.}
\begin{questions}
\item A : $\ce{C2H2}$ (éthyne) ; B : $\ce{CH3Cl}$ (chlorométhane) ; C : $\ce{C2H4}$ (éthène).
\item Les liaisons triple et double correspondent aux tiges triple et double : les modèles sont \res{\text{compatibles}}.
\par\smallskip A : \chemfig{H-C~C-H}\qquad B : \chemfig{H-C(-[2]H)(-[6]H)-\charge{0=\|,90=\|,270=\|}{Cl}}\qquad
C : \chemfig{H-[1]C(-[3]H)=C(-[1]H)-[7]H}
\item A : \res{\text{linéaire}} ; B : \res{\text{tétraédrique}} ; C : \res{\text{plane}}, d'angles voisins de \qty{120}{\degree} autour de chaque carbone.
\item Seule la molécule \res{\text{B}} est polaire : C--Cl est polarisée et n'est pas compensée. A et C ne comportent que des
liaisons C--C et C--H, non polarisées.
\end{questions}
\end{corrige}
\end{blocLivre}
